% !TEX root = ../main.tex

% 中英标题：\chapter{中文标题}[英文标题]
\chapter{预计困难及解决方案}

\section{预计困难与技术难点}
（1）课题涉及的数据量大，且具有高维、高噪声的特性，这使得模型不仅在数据预处理上面临着一些困难，而且在抗噪能力和处理大规模数据的能力上有较高的要求，怎样将蛋白质序列和结构数据转换为统一潜在空间分布也是本研究的一个重难点。

（2）扩散模型的设计包括噪声添加和去噪两个过程，在模拟蛋白质数据的扩散过程时，需要精准地实现高斯噪声的逐步添加，并通过神经网络恢复原始蛋白质表示。优化扩散分类器以实现有效的多标签分类，特别是在低蛋白质序列同源相似性的条件下，能够提供稳定的预测结果，是一个重要的技术挑战。

（3）通过3D点云表征蛋白质结构数据并进行多模态生物数据预训练，与利用扩散模型等生成方法进行分类是近两年才兴起的发展方向，可参考的论文和方法有限，这使得本课题在研究时面临着一些阻力。

(4)在分类类别数量庞大的情况下，计算每个类别的得分通常需要进行多次采样，这将导致计算成本的显著增加。如何设计一种有效的算法和网络框架以降低计算复杂度，已成为该领域目前面临的主要挑战之一。

\section{解决方案}

（1）针对数据集的问题，设计改进的变分自编码器预训练架构，例如引入自注意力机制，增强潜在空间表示的表达能力，确保其能捕捉蛋白质数据的关键特征。

（2）为了提高进展速度和效率，本课题将一边进行实验一边吸收已经提出的蛋白质功能预测领域的相关方法，针对不同算法的问题尝试改进和优化，争取探索最有效的算法来进行研究。

（3）持续阅读大量相关论文，紧跟学术前沿动态，不断尝试最新的多模态预训练方法，拓展科研思路。

